\documentclass[11pt]{article}

\usepackage[margin=1in]{geometry}
\usepackage{amsmath,amssymb,amsthm,mathtools}
\usepackage{microtype}
\usepackage[hidelinks]{hyperref}
\usepackage{xcolor}

\newtheorem{theorem}{Theorem}
\newtheorem{lemma}[theorem]{Lemma}
\newtheorem{corollary}[theorem]{Corollary}
\newcommand{\R}{\mathbb R}
\newcommand{\1}{\mathbf 1}
\newcommand{\dd}{\,d}

\title{A Log-Free $p$-Triangle Inequality for $L^{p,r}$\\
\large A candidate full solution to an open problem in arXiv:1801.10570}
\author{Automated mathematics run \texttt{fa\_banach\_001}, lane 16}
\date{13 August 2026}

\begin{document}
\maketitle

\begin{center}
\fbox{\parbox{0.91\textwidth}{\textbf{Status.}
Candidate full solution, likely valid, awaiting expert review.  The theorem
removes the logarithmic factor whose necessity is asked about by
Seeger--Trebels.  The proof is self-contained and quantitative for the
Lorentz quasi-norm normalization used in the source.}}
\end{center}

\section{Source problem}

Let $(X,\mu)$ be a nonatomic measure space.  For $0<p,r<\infty$, write
\begin{equation}\label{eq:lorentz}
 \|f\|_{p,r}
 =\left(\frac rp\int_0^\infty
       [t^{1/p}f^*(t)]^r\frac{\dd t}{t}\right)^{1/r},
 \qquad
 \|f\|_{p,\infty}=\sup_{t>0}t^{1/p}f^*(t).
\end{equation}
Here $f^*$ is the nonincreasing rearrangement of $|f|$.

Seeger and Trebels prove that, when $0<p<1$ and $p<r<\infty$,
\begin{equation}\label{eq:ptriangle}
 \left\|\sum_k f_k\right\|_{p,r}
 \le C(p,r)\left(\sum_k\|f_k\|_{p,r}^p\right)^{1/p}.
\end{equation}
Their bound for $C(p,r)$ contains
\[
 \left(1+\frac pr\log\frac1{1-p}\right)^{1/p-1/r}.
\]
They ask explicitly whether this logarithmic term is necessary
\cite[Open problem after (13), and Appendix B]{ST}.  We give a negative
answer.

\section{Idea of the proof}

Put $\theta=1-p/r$, so that $r=p/(1-\theta)$.  The two endpoints behind
the result are transparent.  At $r=p$, inequality \eqref{eq:ptriangle} is
the pointwise $p$-triangle inequality in $L^p$ and has constant one.  At
$r=\infty$, a common amplitude cutoff gives a weak-$L^p$ estimate of size
$(1-p)^{-1/p}$.

The main point is to interpolate these endpoints without invoking an
abstract quasi-Banach interpolation theorem or losing control of constants.
For every simple $f$, nonatomicity supplies a rank coordinate $\rho$ with
$|f|=f^*(\rho)$.  The exact factorization
\[
 |f|=u^{1-\theta}v^\theta,\qquad
 v=\rho^{-1/p},\qquad
 u=|f|^{1/(1-\theta)}
       \rho^{\theta/[p(1-\theta)]}
\]
places $u$ in $L^p$ and $v$ in weak $L^p$.  Its normalization is explicit:
\[
 \|u\|_p^p=(1-\theta)\|f\|_{p,r}^r,
 \qquad \|v\|_{p,\infty}\le1.
\]
After rescaling each pair of factors, discrete H\"older turns a sum of
products into a product of sums.  A two-threshold rearrangement inequality
then returns that product to $L^{p,r}$.  Optimizing the two thresholds
produces only the binary-entropy factor
$\theta^{-\theta}(1-\theta)^{-(1-\theta)}\le2$, not a logarithm.

\section{The log-free theorem}

For $0<\theta<1$, set
\begin{equation}\label{eq:entropy}
 \mathcal H(\theta)
 =\theta^{-\theta}(1-\theta)^{-(1-\theta)}.
\end{equation}
As usual, the endpoint values are interpreted by continuity, and
$1\le\mathcal H(\theta)\le2$.

\begin{theorem}[Entropy-refined log-free bound]\label{thm:main}
Let $0<p<1$ and $p<r<\infty$, and put $\theta=1-p/r$.  On every nonatomic
measure space, every finite or countable family with
$\sum_k\|f_k\|_{p,r}^p<\infty$ satisfies
\begin{equation}\label{eq:mainsharp}
 \left\|\sum_k f_k\right\|_{p,r}
 \le
 \mathcal H(\theta)^{1/p}
 p^{-\theta}(1-p)^{-\theta/p}
 \left(\sum_k\|f_k\|_{p,r}^p\right)^{1/p}.
\end{equation}
Consequently, with an absolute constant $A=2e^{1/e}$,
\begin{equation}\label{eq:clean}
 C(p,r)\le A^{1/p}
 \left(\frac1{1-p}\right)^{1/p-1/r}.
\end{equation}
In particular, the logarithmic factor in the Seeger--Trebels bound is not
necessary.
\end{theorem}

The rest of the packet proves the theorem from three elementary lemmas.

\section{Three lemmas}

\subsection{A self-contained weak endpoint}

\begin{lemma}[Weak $p$-summation]\label{lem:weak}
For $0<p<1$ and every finite family of measurable functions,
\begin{equation}\label{eq:weak}
 \left\|\sum_k |g_k|\right\|_{p,\infty}
 \le D_p\left(\sum_k\|g_k\|_{p,\infty}^p\right)^{1/p},
 \qquad
 D_p=p^{-1}(1-p)^{-1/p}.
\end{equation}
\end{lemma}

\begin{proof}
Put $a_k=\|g_k\|_{p,\infty}$ and
$S^p=\sum_k a_k^p$.  Fix $\alpha>0$ and a common cutoff $\beta>0$.
Let
\[
 E_\beta=\bigcup_k\{|g_k|>\beta\},
 \qquad h_k=|g_k|\1_{\{|g_k|\le\beta\}}.
\]
The definition of the weak norm gives
\begin{equation}\label{eq:E}
 \mu(E_\beta)\le \beta^{-p}S^p.
\end{equation}
Layer cake and the same weak estimate give
\begin{equation}\label{eq:lowL1}
 \int_X h_k\dd\mu
 \le\int_0^\beta\mu(|g_k|>s)\dd s
 \le \frac{a_k^p\beta^{1-p}}{1-p}.
\end{equation}
Outside $E_\beta$, the sums of the $|g_k|$ and the $h_k$ agree.  Hence
Chebyshev's inequality, \eqref{eq:E}, and \eqref{eq:lowL1} imply
\[
 \mu\left(\sum_k|g_k|>\alpha\right)
 \le S^p\left(\beta^{-p}
       +\frac{\beta^{1-p}}{(1-p)\alpha}\right).
\]
The choice $\beta=p\alpha$ minimizes the displayed expression and makes it
\[
 S^p\alpha^{-p}\frac{p^{-p}}{1-p}.
\]
Taking the supremum of
$\alpha\,\mu(\sum_k|g_k|>\alpha)^{1/p}$ proves \eqref{eq:weak}.
\end{proof}

\subsection{Quantitative rank factorization}

\begin{lemma}[Strong--weak factorization]\label{lem:factor}
Let $0<p<1$, $0<\theta<1$, and $r=p/(1-\theta)$.  If $f$ is a
nonnegative simple function with support of finite measure on a nonatomic
measure space, then there are $u,v\ge0$ such that
\begin{equation}\label{eq:factor}
 f=u^{1-\theta}v^\theta,\qquad
 \|v\|_{p,\infty}\le1,\qquad
 \|u\|_p^p=(1-\theta)\|f\|_{p,r}^r.
\end{equation}
\end{lemma}

\begin{proof}
Write the positive values of $f$ in decreasing order as
$c_1>\cdots>c_m>0$ and let $E_j=\{f=c_j\}$.  Set
$M_j=\sum_{i\le j}\mu(E_i)$ and $M_0=0$.  Nonatomicity, by recursive
bisection, supplies on every $E_j$ a measurable variable uniformly
distributed on $(M_{j-1},M_j)$.  Combining these variables gives a rank
map
\[
 \rho:\operatorname{supp}f\longrightarrow(0,M_m)
\]
whose pushforward is Lebesgue measure and for which
$f=f^*(\rho)$ almost everywhere.  This also handles flat levels: each level
is spread across its full rearrangement interval.

Define
\[
 v=\rho^{-1/p},
 \qquad
 u=f^{1/(1-\theta)}
       \rho^{\theta/[p(1-\theta)]}.
\]
The factorization in \eqref{eq:factor} is immediate.  Since $\rho$ is
uniform,
\[
 \mu(v>\lambda)=\min\{M_m,\lambda^{-p}\},
\]
and therefore $\|v\|_{p,\infty}\le1$.  Also, using
$r=p/(1-\theta)$ and the normalization \eqref{eq:lorentz},
\begin{align*}
 \|u\|_p^p
 &=\int_0^{M_m}f^*(t)^{p/(1-\theta)}
       t^{\theta/(1-\theta)}\dd t\\
 &=\int_0^\infty f^*(t)^r t^{r/p-1}\dd t
  =(1-\theta)\|f\|_{p,r}^r.
\end{align*}
This proves all assertions.
\end{proof}

\subsection{Returning a product to the Lorentz space}

\begin{lemma}[Product rearrangement]\label{lem:product}
Let $0<p<\infty$, $0<\theta<1$, and $r=p/(1-\theta)$.  Then
\begin{equation}\label{eq:product}
 \|u^{1-\theta}v^\theta\|_{p,r}
 \le K_\theta\|u\|_p^{1-\theta}
                 \|v\|_{p,\infty}^\theta,
\end{equation}
where
\begin{equation}\label{eq:K}
 K_\theta=(1-\theta)^{-2/r}\theta^{-\theta/p}.
\end{equation}
\end{lemma}

\begin{proof}
For every $0<a<1$ and $t>0$, the elementary two-threshold rearrangement
bound is
\begin{equation}\label{eq:two-threshold}
 (u^{1-\theta}v^\theta)^*(t)
 \le u^*(at)^{1-\theta}v^*((1-a)t)^\theta.
\end{equation}
Indeed, away from the union of the sets where $u$ and $v$ exceed the two
right-hand thresholds, whose total measure is at most $t$, the product is
at most the right-hand side.  The weak estimate for $v$ and
$r(1-\theta)=p$ now give
\begin{align*}
 \|u^{1-\theta}v^\theta\|_{p,r}^r
 &\le \frac rp(1-a)^{-\theta r/p}
       \|v\|_{p,\infty}^{\theta r}
       \int_0^\infty u^*(at)^p\dd t\\
 &=\frac{r}{pa}(1-a)^{-\theta r/p}
       \|u\|_p^p\|v\|_{p,\infty}^{\theta r}.
\end{align*}
Thus \eqref{eq:product} holds with
\[
 \left(\frac{r}{pa}\right)^{1/r}(1-a)^{-\theta/p}.
\]
Choosing $a=1-\theta=p/r$ yields \eqref{eq:K}.
\end{proof}

\section{Proof of Theorem~\ref{thm:main}}

We first take finitely many nonnegative simple functions of finite-measure
support.  Put
\[
 F_k=\|f_k\|_{p,r},
 \qquad
 S=\left(\sum_kF_k^p\right)^{1/p}.
\]
Ignore zero terms.  Lemma~\ref{lem:factor} gives
$f_k=u_{k,0}^{1-\theta}v_{k,0}^\theta$ and
\begin{equation}\label{eq:rawbalance}
 \|u_{k,0}\|_p^{1-\theta}
 \|v_{k,0}\|_{p,\infty}^\theta
 \le (1-\theta)^{1/r}F_k.
\end{equation}
Rescale the two factors reciprocally, preserving their product, so that the
new factors $u_k,v_k$ obey
\begin{equation}\label{eq:scaled}
 f_k=u_k^{1-\theta}v_k^\theta,\qquad
 \|u_k\|_p=F_k,\qquad
 \|v_k\|_{p,\infty}
 \le (1-\theta)^{1/(r\theta)}F_k.
\end{equation}
For completeness, if $u_k=c_ku_{k,0}$ with
$c_k=F_k/\|u_{k,0}\|_p$, take
$v_k=c_k^{-(1-\theta)/\theta}v_{k,0}$; the final estimate in
\eqref{eq:scaled} is exactly \eqref{eq:rawbalance} raised to the power
$1/\theta$.

Let $U=\sum_k u_k$ and $V=\sum_kv_k$.  Discrete H\"older's inequality at
each point gives
\begin{equation}\label{eq:discrete-holder}
 \sum_k f_k
 \le U^{1-\theta}V^\theta.
\end{equation}
The scalar inequality $(\sum a_k)^p\le\sum a_k^p$ and
\eqref{eq:scaled} give
\begin{equation}\label{eq:U}
 \|U\|_p\le S.
\end{equation}
Lemma~\ref{lem:weak} and \eqref{eq:scaled} give
\begin{equation}\label{eq:V}
 \|V\|_{p,\infty}
 \le D_p(1-\theta)^{1/(r\theta)}S.
\end{equation}
Applying Lemma~\ref{lem:product} to \eqref{eq:discrete-holder}, then using
\eqref{eq:U} and \eqref{eq:V}, yields
\begin{align*}
 \left\|\sum_k f_k\right\|_{p,r}
 &\le K_\theta D_p^\theta(1-\theta)^{1/r}S\\
 &=\mathcal H(\theta)^{1/p}
   p^{-\theta}(1-p)^{-\theta/p}S.
\end{align*}
The last identity uses $1/r=(1-\theta)/p$ and
$D_p=p^{-1}(1-p)^{-1/p}$.  This is \eqref{eq:mainsharp}.

For arbitrary functions, first use
$|\sum_k f_k|\le\sum_k|f_k|$.  Approximate each $|f_k|$ increasingly by
finite-support simple functions.  Such an approximation exists because
$\mu(|f_k|>\varepsilon)<\infty$ for every $\varepsilon>0$.  The Fatou
property of the Lorentz quasi-norm passes the proved estimate to the limit.
Countable families follow by applying the finite result to partial sums and
using Fatou once more.

Finally, $\theta/p=1/p-1/r$, $\mathcal H(\theta)\le2$, and
\[
 p^{-p\theta}\le p^{-p}\le e^{1/e}\qquad(0<p<1).
\]
Raising the refined constant in \eqref{eq:mainsharp} to the $p$th power
therefore gives \eqref{eq:clean} with $A=2e^{1/e}$.  The logarithm is gone.
\qed

\section{Endpoint checks and scope}

The refined bound has the correct qualitative interpolation limits.  As
$r\downarrow p$, one has $\theta\downarrow0$ and its right-hand constant
tends to one, exactly matching $L^{p,p}=L^p$.  As $r\to\infty$, one has
$\theta\uparrow1$ and the constant tends to
$D_p=p^{-1}(1-p)^{-1/p}$, the weak endpoint proved above.

The theorem answers the source's stated yes/no issue: no logarithmic factor
is required in a uniform upper bound.  It does not claim that the displayed
entropy-refined constant is the exact best value of $C(p,r)$.  Determining
the exact optimum remains a separate extremal problem.

\begin{thebibliography}{9}

\bibitem{ST}
A.~Seeger and W.~Trebels,
\emph{Embeddings for spaces of Lorentz--Sobolev type},
Math. Ann. \textbf{373} (2019), 1017--1056;
arXiv:1801.10570.

\bibitem{STW}
E.~M. Stein, M.~H. Taibleson, and G.~Weiss,
\emph{Weak type estimates for maximal operators on certain $H^p$ classes},
Rend. Circ. Mat. Palermo (2), Suppl. 1 (1981), 81--97.

\end{thebibliography}

\end{document}
